\chapter{Best-First Search and Heuristic Guarantees}
\label{sec:p12-best-first-search-and-heuristic-guarantees}
\index{search!best-first}
\index{search!A*}
\index{search!A**}
\index{heuristic!admissibility}
\index{heuristic!consistency}
\index{search!reopening}
\index{search!computational optimality}
\index{search!global compatibility}

A heuristic can guide the search for a cheap solution without proving that
the first solution found is cheapest. A lower bound supplies the missing
certificate, provided the search preserves the paths needed by that
certificate. A separate question concerns computational efficiency: which
other algorithms must inspect the same states? The two questions have
different hypotheses. This treatment follows the distinctions developed by
Dechter and Pearl \cite{DechterPearl1985BestFirst}.

\section{Search objects and the information interface}

\begin{hbdefinition}{def:p12-hsearch-instance}{A weighted state graph}
An instance is \(I=(G,s,\Gamma,c,h)\). Here \(G=(V,E)\) is a directed
graph, \(s\in V\) is the initial state, \(\Gamma\subseteq V\) is the set
of goals, and each edge has a finite nonnegative real cost \(c(u,v)\).
A path record \(P=(v_0,\ldots,v_k)\) includes its edges, starts at
\(v_0=s\), and has cost
\[
 g(P)=\sum_{i=0}^{k-1}c(v_i,v_{i+1}).
\]
The empty-edge path at \(s\) has cost zero. A solution ends at a goal.
The notation \(P_{s,u}\) denotes the prefix ending at the specified
occurrence of \(u\); segments of a record are indexed by occurrences
when a state is repeated.
Write \(g^*(v)\) for the infimum of costs from \(s\) to \(v\), and
\(h^*(v)\) for the infimum of costs from \(v\) to a goal; an empty
set of possible paths has infimum \(+\infty\). Write \(C^*=h^*(s)\).
When an optimal path is used, attainment of this infimum is required.

The heuristic \(h:V\to[0,\infty)\) is finite at every state and is
zero at goals. The graph is \emph{locally finite} when each state has
finitely many successors. Expanding a non-goal state reveals its complete
successor list and edge costs. Generation reveals a successor's identity,
goal flag and heuristic. State equality is decidable in this interface.
For oracle-expansion comparisons, neither the complete graph, its exact
vertex count nor a complete list of goals is supplied in advance.
Goals are terminal in this search interface: their successor lists are
not queried. With nonnegative additive costs, truncating a path at its
first goal cannot make the solution more expensive.
\end{hbdefinition}

A node of the path-unfolding tree is a path record, whereas a vertex of
\(G\) is a state. Several records can end at the same state. Counting
states, counting path records, counting repeated expansions and measuring
elapsed computation therefore give different quantities.
\begin{hbdefinition}{def:p12-hsearch-expansion-count}{The compared expansion set}
For the comparisons below, \(E_A(I)\) denotes the set of distinct non-goal states
expanded by algorithm \(A\). Selecting a goal and returning without
generating its successors does not expand it in this convention.
\end{hbdefinition}

\begin{hbdefinition}{def:p12-hsearch-heuristics}{Admissibility and consistency}
The heuristic is \emph{admissible} when \(h(v)\le h^*(v)\) for every
state. It is \emph{consistent} when, for every edge,
\[
 h(u)\le c(u,v)+h(v).
\]
The latter condition concerns pairs of states, not just the accuracy of
individual estimates. The normalization \(h(t)=0\) at every goal is part
of both the algorithmic and comparison interface.
\end{hbdefinition}

\begin{hbtheorem}{thm:p12-hsearch-consistency-bound}{proved=annex:p12.hsearch-consistency-bound}{Consistency supplies a lower bound}
A consistent normalized heuristic is admissible. Along every path record,
the values \(g+h\) are nondecreasing.
\end{hbtheorem}

\section{An executable A* graph procedure}

\begin{hbmethod}{def:p12-hsearch-algorithm}{A* with duplicate detection and reopening}
Store a label \(g(v)\), initially \(+\infty\), and a witnessing path
record \(W(v)\) for every reached state. Set \(g(s)=0\), \(W(s)=(s)\)
and \(\mathrm{OPEN}=\{s\}\), with \(\mathrm{CLOSED}=\varnothing\).
Repeat the following steps.
\begin{enumerate}
\item If OPEN is empty, return failure.
\item Select a state \(u\) in OPEN with minimum
  \(f(u)=g(u)+h(u)\). Among equal scores prefer goals; remaining ties
  use a declared rule based on the revealed information. Remove \(u\)
  from OPEN.
\item If \(u\) is a goal, return \(W(u)\). In particular, do not
  return merely because a goal was generated.
\item Put \(u\) in CLOSED and reveal all successors. For each
  \((u,v)\), compute \(q=g(u)+c(u,v)\). If \(q<g(v)\), replace
  \(g(v)\) by \(q\), set \(W(v)\) to the extension of the current
  record \(W(u)\) by \((u,v)\), remove \(v\) from CLOSED if present,
  and put \(v\) in OPEN. If \(q\ge g(v)\), retain the old label.
\end{enumerate}
Path records are persistent: changing a state's current label does not
silently change the cost or ancestry stored in an older record. A heap
implementation may contain obsolete entries, but must discard an entry
whose version is no longer the state's current OPEN label. Re-expansion
after a strict improvement is a real operation even if the successor list
has already been cached.
\end{hbmethod}

Since \(h(v)\) is fixed for a given state, comparing two tentative
\(f\)-values at that state is exactly comparing their \(g\)-values.
This is the duplicate-pruning property used by A*. An arbitrary
path-dependent score need not have it. For such a score, keep distinct
path records unless the discarded record is proved no better under
every possible common continuation.

\begin{hblemma}{thm:p12-hsearch-frontier}{proved=annex:p12.hsearch-frontier}{A frontier witness on a fixed path}
Let \(P\) be a finite path from \(s\) to a goal. Until A* terminates,
some state \(v\) on \(P\) is in OPEN with
\(g(v)\le g(P_{s,v})\), unless a goal of no larger path cost has
already been selected. For an optimal path and admissible \(h\), this
witness satisfies \(f(v)\le C^*\).
\end{hblemma}

\begin{hbtheorem}{thm:p12-hsearch-finite}{proved=annex:p12.hsearch-finite}{Finite-graph termination and failure}
On a finite graph with nonnegative edge costs, the stated A* procedure
terminates for every finite heuristic satisfying the goal normalization.
It returns failure exactly when no goal is reachable from \(s\).
This termination statement does not require heuristic admissibility.
\end{hbtheorem}

\begin{hbtheorem}{thm:p12-hsearch-optimal}{proved=annex:p12.hsearch-optimal}{Least-cost solution}
If an optimal path exists, \(h\) is admissible, and the stated A*
procedure terminates by selecting a goal, the returned path has cost
\(C^*\). Before termination every selected non-goal state satisfies
\(g(v)+h(v)\le C^*\). Thus finite-graph termination gives a complete
least-cost search procedure on finite graphs.
\end{hbtheorem}

\begin{hbtheorem}{thm:p12-hsearch-infinite}{proved=annex:p12.hsearch-infinite}{A sufficient condition on infinite graphs}
Suppose \(G\) is locally finite, a goal is reachable, and some constant
\(\delta>0\) bounds every edge cost below. With finite nonnegative
\(h\), A* terminates with a solution. If \(h\) is also admissible,
the solution is cheapest. Neither assertion promises termination with
failure on an infinite unsolvable graph.
\end{hbtheorem}

\begin{hbtheorem}{thm:p12-hsearch-no-reopen}{proved=annex:p12.hsearch-no-reopen}{When CLOSED really is final}
With a consistent heuristic, every state selected by A* has its least
start-to-state cost \(g(v)=g^*(v)\). Consequently no CLOSED state needs
reopening. Admissibility alone does not justify that implementation.
\end{hbtheorem}

\section{Deriving and using a heuristic}

\begin{hbdefinition}{def:p12-hsearch-relaxation}{An edge-preserving relaxation}
Let \(\rho\) map original states to states of another weighted graph.
Every original goal maps to a relaxed goal. For each original edge
\((u,v)\), the relaxed problem has a path from \(\rho(u)\) to
\(\rho(v)\) costing at most \(c(u,v)\). Let \(d_R(x)\) be the exact
least cost from relaxed state \(x\) to a relaxed goal, and suppose
\(d_R(\rho(v))\) is finite. Define \(h_R(v)=d_R(\rho(v))\).
\end{hbdefinition}

\begin{hbtheorem}{thm:p12-hsearch-relaxation}{proved=annex:p12.hsearch-relaxation}{The relaxation certificate}
The heuristic \(h_R\) is admissible and consistent. Its guarantee uses
the common edge-preserving relaxation and exact relaxed optimum; an
unrelated convenient score or an arbitrary suboptimal relaxed solution
cost does not supply this certificate.
\end{hbtheorem}

\begin{hbexample}{def:p12-hsearch-trace}{A complete frontier trace}
The only edges are \(s\to a\) of cost 2, \(a\to t\) of cost 2,
\(s\to b\) of cost 1 and \(b\to t\) of cost 10, with sole goal
\(t\). Take \(h(s)=0\), \(h(a)=2\), \(h(b)=10\), \(h(t)=0\).
After expanding \(s\), OPEN contains
\((a:g=2,h=2,f=4)\) and \((b:g=1,h=10,f=11)\).
Select \(a\), generate \((t:g=4,h=0,f=4)\), then select \(t\)
and return \(s,a,t\), cost 4. The parent of \(t\) is \(a\);
\(b\) remains unexpanded. For a concrete relaxation producing this
entire heuristic, retain every original edge and add a zero-cost edge
\(s\to t\). The relaxed remaining costs are exactly
\(0,2,10,0\) at \(s,a,b,t\); each original edge is still available
at its original cost.

If instead \(h(a)=20\), OPEN after \(s\) has scores 22 and 11.
Expand \(b\), generate \(t\) with \(g=f=11\), and select it before
\(a\). The returned cost is 11, although the optimum is 4. This is a
solution-quality failure caused by overestimation.
\end{hbexample}

\begin{hbexample}{def:p12-hsearch-generated-goal}{Generation is not the stopping certificate}
With edges \(s\to t:10\), \(s\to a:1\), \(a\to t:1\)
and \(h=0\), expanding \(s\) generates a goal of cost 10.
Its OPEN score is nevertheless larger than \(a\)'s score 1.
Selecting \(a\) first improves the goal label to 2. Returning on
generation would lose even this elementary uniform-cost guarantee.
\end{hbexample}

\begin{hbexample}{def:p12-hsearch-reopening}{An admissible inconsistent estimate}
Use edges \(s\to a:3\), \(s\to b:1\), \(b\to a:1\),
\(a\to t:2\), with sole goal \(t\), and set
\(h(s)=h(a)=h(t)=0\), \(h(b)=3\).
All estimates are admissible, but \(h(b)>1+h(a)\).
A* first expands \(a\) with \(g=3\) and generates \(t\) at cost 5.
It next expands \(b\), obtaining a better label \(g(a)=2\).
Reopening \(a\) improves \(t\) to cost 4 and changes its witnessing
path to \(s,b,a,t\). A procedure that permanently discards improvements
to CLOSED returns cost 5 on this same input.
\end{hbexample}

\begin{hbexample}{def:p12-hsearch-termination-failures}{Two infinite-search boundaries}
Local finiteness alone is insufficient. Let \(s\) have a goal successor
of cost 2 and a chain whose successive edge costs are
\(1/2,1/4,1/8,\ldots\). With \(h=0\), every chain prefix costs less
than 1 and is selected before the goal forever. This example has strictly
positive costs but no positive uniform lower bound.

Alternatively, on an infinite unit-cost chain with no goal, A* with
\(h=0\) never empties OPEN. Existence of a solution is an essential
premise of Theorem~\ref{thm:p12-hsearch-infinite}.
\end{hbexample}

\section{What an efficiency comparison means}

\begin{hbdefinition}{def:p12-hsearch-comparison}{Equal information and four comparison tests}
An equal-information algorithm begins at \(s\), expands only generated
states, receives the same successor, goal and heuristic information, and
has no oracle for ungenerated portions of the graph. It may process known
information differently. In particular, learning a tighter bound from an
already revealed edge does not constitute extra external information.
The comparison excludes a backward search started from an undisclosed goal
or additional independently seeded search trees.
The instance domain here consists of abstract weighted graphs. Extra known
restrictions, such as unit edge costs or a fixed geometric state model,
define a different comparison domain: the adversarial completions used
below must belong to that domain before its lower bounds can be asserted.

Fix a domain \(\mathcal I\), a competitor class \(\mathcal A\), and
a family \(\{A_\tau\}\) of goal-preferring A* tie policies. The four
tests distinguish the following quantifiers:
\begin{enumerate}
\item Universal expansion containment:
 \(\forall I\,\forall\tau\,\forall B\in\mathcal A,
 E_{A_\tau}(I)\subseteq E_B(I)\).
\item Containment up to a competitor-dependent choice of ties:
 \(\forall I\,\forall B\in\mathcal A\,\exists\tau,
 E_{A_\tau}(I)\subseteq E_B(I)\).
\item Pointwise absence of a proper-subset competitor, for a specified
 policy: \(\forall I\,\forall B\in\mathcal A,
 E_B(I)\not\subsetneq E_{A_\tau}(I)\).
\item Absence of strict uniform domination of a specified policy:
 if \(E_{A_\tau}(I_1)\setminus E_B(I_1)\ne\varnothing\) for some
 \(I_1\), then
 \(E_B(I_2)\setminus E_{A_\tau}(I_2)\ne\varnothing\) for some
 \(I_2\).
\end{enumerate}
These are the distinctions underlying Dechter--Pearl types 0--3.
The second test does not require one tie policy to work simultaneously
against every competitor. Tests involving a specified policy must not be
silently quantified over all policies when comparing their strengths.
\end{hbdefinition}

For example, two equal-cost solution branches with perfect estimates let
different algorithms choose different branches. A tie choice reproducing
one competitor's branch need not reproduce another's. A gratuitous
expansion made by a different tie policy cannot be charged to the first
policy. This is why the quantifiers belong in the statement of a claimed
efficiency theorem.

\begin{hbdefinition}{def:p12-hsearch-domains}{Admissible and consistent comparison domains}
Let \(\mathcal I_{\rm adm}\) contain solvable finite instances with
admissible estimates, and let \(\mathcal I_{\rm con}\) impose
consistency as well. Let \(\mathcal A_{\rm adm}\) be the algorithms
that return cheapest solutions on every instance in
\(\mathcal I_{\rm adm}\). The guarantee is across that entire class,
not just correctness on the particular graph being measured.
\end{hbdefinition}

\begin{hbdefinition}{def:p12-hsearch-sure-set}{A strict contour and a nonpathological instance}
For an admissible instance, put
\[
 S(I)=\{v\notin\Gamma:\text{some simple path }P:s\leadsto v
 \text{ has }g(P_{s,u})+h(u)<C^*
 \text{ at every noninitial }u\in P\}.
\]
The initial state is included when it is not a goal. An instance is
nonpathological here when it has an optimal path on which
\(h(u)<h^*(u)\) at every non-goal state, including \(s\).
\end{hbdefinition}

\begin{hbtheorem}{thm:p12-hsearch-sure-set}{proved=annex:p12.hsearch-sure-set}{Strict contours are expanded}
A* expands every state in \(S(I)\). On a nonpathological finite
instance it expands exactly \(S(I)\), independently of remaining ties.
\end{hbtheorem}

\begin{hbtheorem}{thm:p12-hsearch-consistent-efficiency}{proved=annex:p12.hsearch-consistent-efficiency}{The consistent equal-information lower bound}
On a consistent finite instance, every algorithm in
\(\mathcal A_{\rm adm}\) expands every state in \(S(I)\).
Consequently A* has universal expansion containment on nonpathological
consistent instances, and competitor-dependent containment up to ties on
all consistent finite instances.
\end{hbtheorem}

\begin{hbexample}{def:p12-hsearch-information-trade}{An admissible estimate does not imply minimum expansions}
Use \(s\to a:3\), \(s\to b:1\), \(s\to c:1\),
\(b\to a:1\), \(c\to t:10\), with sole goal \(t\) and
\(a\) a dead end. Set \(h(s)=0,h(a)=1,h(b)=20,h(c)=9,h(t)=0\).
The infinite true remaining costs at \(a,b\) make these finite
estimates admissible. A* expands \(a\) at score 4, then \(c\) at
score 10, and returns cost 11 without expanding \(b\).

An algorithm can instead inspect \(c\), discover a solution of cost
11, and then inspect \(b\). The revealed edge \(b\to a\), together
with admissibility, implies
\(h^*(a)\ge h(b)-1=19\). Every still undiscovered solution must pass
through \(a\), with prefix cost at least 2, so none can improve 11.
The algorithm may stop without expanding \(a\). This saves \(a\)
but pays for \(b\). It exemplifies the information trade in the
paper's Figure 6(a); it is not a failure of A*'s returned cost.
\end{hbexample}

\begin{hbexample}{def:p12-hsearch-pathological-subset}{A proper-subset competitor}
Use \(s\to a:3\), \(s\to b:1\), \(b\to a:1\),
\(b\to t:10\), with \(a\) a dead end, and
\(h(s)=h(t)=0,h(a)=1,h(b)=10\). A* must expand \(a\) before
\(b\), then returns the goal at cost 11. A competitor that inspects
\(b\) first learns \(h^*(a)\ge 9\), so all unseen solutions through
\(a\) cost at least \(2+9=11\). It can return its known solution
without expanding \(a\): its expansion set is a proper subset.
The optimal path has \(h(b)=h^*(b)\), so the instance is pathological.
Both competitors described here can be made admissible on all finite
inputs by using only the justified stopping certificate and otherwise
continuing exhaustive search.
\end{hbexample}

\begin{hbexample}{def:p12-hsearch-subset-cost-boundary}{The prefix cost in a pruning certificate}
In the preceding pattern, change only the direct edge \(s\to a\)
from cost 3 to cost 1. The same revealed information now bounds an
unseen solution through \(a\) only by \(1+9=10\), which does
not justify returning the known cost 11. A hidden goal continuation
from \(a\) of cost 9 gives an optimum of 10, while preserving
\(h(a)=1\) and \(h(b)=10\) as admissible estimates. Thus a
pruning certificate must use the cheapest available prefix, not merely
the prefix that passes through the state supplying its heuristic bound.
\end{hbexample}

\begin{hbtheorem}{thm:p12-hsearch-tradeoff}{proved=annex:p12.hsearch-tradeoff}{The nonpathological expansion trade}
On a finite nonpathological admissible instance, if an algorithm in
\(\mathcal A_{\rm adm}\) skips a state in \(S(I)\), it expands
some state outside \(S(I)\). Thus no admissible competitor has a
proper-subset expansion set on such an instance. There is nevertheless
no single admissible algorithm whose expansion set is contained in every
admissible competitor's on every admissible finite instance.
\end{hbtheorem}

\begin{hbdefinition}{def:p12-hsearch-cost-compatibility}{Costwise compatibility outside the admissible domain}
Call \(B\) costwise compatible with the A* family when, on every
finite solvable nonnegative-cost instance with finite normalized
\(h\), it returns a path costing no more than the path returned by
any goal-preferring A* tie policy. This condition includes cases in which
\(h\) overestimates. It is stronger than merely requiring \(B\)
to be optimal whenever A* happens to be optimal.
\end{hbdefinition}

\begin{hbtheorem}{thm:p12-hsearch-cost-compatibility}{proved=annex:p12.hsearch-cost-compatibility}{Expansion containment for a costwise competitor}
On every admissible finite instance, a costwise-compatible algorithm
expands each state that some A* run expands at score strictly below
\(C^*\). A* therefore has competitor-dependent containment up to
ties against this class on the entire admissible finite domain, and
universal containment on its nonpathological subdomain.
\end{hbtheorem}

\section{Path barriers and the A** distinction}

\begin{hbdefinition}{def:p12-hsearch-pathmax}{A path-dependent barrier}
For a path record \(P\), define
\[
 F(P)=\max_{u\in P}\bigl(g(P_{s,u})+h(u)\bigr).
\]
The A** procedure considered here uses the same strict \(g\)-improvement
and reopening rule as Definition~\ref{def:p12-hsearch-algorithm}, but
selects a record minimizing \(F\), with goals preferred at equal
\(F\). Store the path record and its barrier together. Its update is
\(F(P\cdot v)=\max(F(P),g(P\cdot v)+h(v))\).
It does not discard a record merely because another record at the same
state has a smaller barrier: duplicate dominance still compares \(g\).
When a better \(g\)-record is accepted, replace the barrier by that
record's newly computed barrier, even if the new barrier is larger.
This separation of the selection score and duplicate criterion is the
essential Appendix B distinction \cite[pp.~534--535]{DechterPearl1985BestFirst}.
\end{hbdefinition}

\begin{hbtheorem}{thm:p12-hsearch-pathmax}{proved=annex:p12.hsearch-pathmax}{Admissibility of the stated A** procedure}
On solvable finite nonnegative-cost graphs with admissible \(h\), A**
terminates with a cheapest solution; on an unsolvable finite graph it
returns failure. The least-cost conclusion also holds on locally finite
solvable graphs whose edge costs have a positive uniform lower bound.
\end{hbtheorem}

\begin{hbtheorem}{thm:p12-hsearch-pathmax-containment}{proved=annex:p12.hsearch-pathmax-containment}{A** containment with matched ties}
On every admissible finite instance, for each goal-preferring A* tie
policy there exists a goal-preferring A** tie policy that expands a
subset of its distinct non-goal states. Some admissible instances let
A** expand a proper subset of the states expanded by every A* tie
policy. The first assertion permits the matching policy to depend on
the instance and the specified A* run; it is not by itself a theorem
of uniform domination by one fixed A** policy.
\end{hbtheorem}

\begin{hbexample}{def:p12-hsearch-barrier-tie}{Why an inherited barrier changes expansions}
Take \(s\to p:1\), \(p\to u:1\), \(p\to t:9\), with
\(u\) a dead end, \(t\) the sole goal, and
\(h(s)=0,h(p)=9,h(u)=h(t)=0\). After expanding \(p\), A* assigns
\(f(u)=2\) and \(f(t)=10\), so must inspect \(u\).
A** assigns both records barrier 10 and selects the goal, returning the
same optimum while skipping \(u\). This is an admissible inconsistent
case, and the saving concerns a boundary tie in the inherited barrier.
\end{hbexample}

\begin{hbexample}{def:p12-hsearch-wrong-duplicate-score}{A barrier is not a duplicate-dominance key}
Take edges \(s\to a:1\), \(s\to b:2\), \(a\to x:4\),
\(b\to x:5\), \(x\to t:5\), with sole goal \(t\), and
\(h(s)=h(t)=0,h(a)=9,h(b)=6,h(x)=1\). These estimates are
admissible. A** first reaches \(x\) through \(b\), at cost 7
and barrier 8, and generates a goal of cost 12. It later expands \(a\)
and obtains a path to \(x\) of cost 5 but barrier 10. Accepting the
lower cost and reopening \(x\) produces the optimum 10. Rejecting
that path because its barrier exceeds 8 leaves the suboptimal goal 12.
The common extension \(x\to t\) thus reverses the preference between
these two prefixes when preference is based only on their barriers.
\end{hbexample}

\section{The restricted path-score competitor class}

\begin{hbdefinition}{def:p12-hsearch-bestfirst-class}{Normalized path-only best-first competitors}
A competitor in \(\mathcal A_{\rm bf}\) selects a smallest-score
OPEN record, with goal preference, using a score \(b(P)\) determined
only by the states, edge costs, heuristic values and goal flags on that
record. In particular it cannot inspect some separately explored branch
when scoring \(P\). At goals require \(b(P)=g(P)\).
It is admissible on all finite instances with admissible \(h\).
If it merges records at a duplicate state using \(b\), require
continuation-safe order preservation:
\[
 b(P)\le b(Q)\ \Longrightarrow\ b(P\cdot R)\le b(Q\cdot R)
\]
for every common continuation \(R\); a strict improvement reopens
a processed state. Otherwise the competitor must keep the distinct
path records needed by its correctness argument. Goal normalization is
a real restriction on the numerical score, not just an ordering convention.
\end{hbdefinition}

\begin{hbtheorem}{thm:p12-hsearch-bestfirst-forcing}{proved=annex:p12.hsearch-bestfirst-forcing}{What an admissible path score cannot conceal}
Fix an admissible finite instance with \(h(s)<C^*\). On every simple path
\(P:s\leadsto v\) whose noninitial prefix scores \(g+h\) are
strictly below \(C^*\), every noninitial score \(b(P_{s,u})\) of
an algorithm in \(\mathcal A_{\rm bf}\) is strictly below \(C^*\).
It follows that the competitor expands every state of \(S(I)\).
Hence A* has universal expansion containment against this class on
nonpathological instances. On all admissible finite instances, A** has
containment against this class up to a competitor-dependent tie choice.
\end{hbtheorem}

\begin{hbtheorem}{thm:p12-hsearch-two-input-score}{proved=annex:p12.hsearch-two-input-score}{A restriction on two-input scores}
If a member of \(\mathcal A_{\rm bf}\) scores non-goal records by
\(b(P)=H(g(P),h(v))\), then
\(H(x,y)\le x+y\) for every finite \(x,y\ge 0\) that its
declared input domain permits. The assertion uses the normalized goal
score and admissibility across that entire domain.
\end{hbtheorem}

\section{Global guarantees and propagation of known costs}
\label{sec:p12-hsearch-cost-closure}

\begin{hbdefinition}{def:p12-hsearch-weak-compatibility}{The weaker optimum-preservation condition}
A weakly compatible algorithm returns an optimum whenever an A* policy
returns an optimum on the same finite solvable instance. It need not
match A*'s returned cost on an instance where A* is suboptimal.
Requiring this implication for every tie policy is still weaker than
costwise compatibility.
\end{hbdefinition}

\begin{hbtheorem}{thm:p12-hsearch-weak-compatibility}{proved=annex:p12.hsearch-weak-compatibility}{Optimum preservation does not force the strict contour}
Under the stated graph-information interface, there is a weakly
compatible algorithm that, on a nonpathological admissible instance,
skips a state in \(S(I)\). Thus the costwise containment theorem
cannot be extended to this weaker premise.
\end{hbtheorem}

\begin{hbtheorem}{thm:p12-hsearch-pending-domination}{proved=annex:p12.hsearch-pending-domination}{Pending label propagation can defeat type-3 optimality}
For the pending-reopening A** procedure of
Definition~\ref{def:p12-hsearch-pathmax}, there is a legal
goal-preferring tie policy and an admissible equal-information algorithm
that uniformly expands a subset of its states and a proper subset on
some finite positive-cost instance. Consequently arbitrary tie policies
of that operational procedure do not satisfy the type-3 test.
\end{hbtheorem}

\begin{hbmethod}{def:p12-hsearch-saturated}{A path-barrier search with complete known-edge relaxation}
Write \(A^\dagger\) for the following explicitly strengthened
procedure. Maintain the finite revealed graph \(R\), the set \(D\)
of non-goals whose successor lists have been revealed, and the frontier
\(U\) of generated states outside \(D\). Initially \(R\) has
the single vertex \(s\) and no edges, \(D=\varnothing\),
\(U=\{s\}\), \(g(s)=0\), \(W(s)=(s)\) and
\(F(W(s))=h(s)\).

Before selecting a new state, propagate every strict \(g\)-improvement
through every already revealed outgoing edge until no improvement is
possible. Keep persistent witnesses and the path barriers \(F\) with
the same update formula as A**. This relaxation processes cached lists;
it does not reveal a new state. Equal-cost alternatives retain the
existing record. Now select a minimum-barrier state in \(U\), with
goals preferred and otherwise the specified tie policy. Return its
witness if it is a goal. Otherwise reveal its full successor list,
place it in \(D\), initialize or improve successor labels, and repeat.
Empty \(U\) means failure.

At selection times its \(g\)-labels are shortest-path costs in
\(R\). This is stronger than updating only descendants in the current
pointer tree: a newly improved path may use a previously revealed edge
that is not a current parent edge. The extra relaxation is not free in
time or memory; it is free only in the distinct-successor-list metric.
\end{hbmethod}

\begin{hblemma}{thm:p12-hsearch-saturated-invariant}{proved=annex:p12.hsearch-saturated-invariant}{The exposed-graph certificate}
On solvable finite instances with admissible heuristics, \(A^\dagger\)
terminates with an optimum. Its successive selected barriers are
nondecreasing. Immediately
before a non-goal \(j\) is selected with barrier \(m\), let
\(d_R\) be distance in the revealed graph. Then
\[
 g(v)=d_R(s,v),\qquad
 d_R(s,x)+h(x)\le m\quad(x\in D),\qquad
 d_R(s,j)+h(j)\le m.
\]
Every generated goal has known least revealed-path cost strictly greater
than \(m\). The same procedure is complete on solvable locally finite
graphs with a positive uniform lower bound on edge costs.
\end{hblemma}

\begin{hbtheorem}{thm:p12-hsearch-saturated-type-three}{proved=annex:p12.hsearch-saturated-type-three}{Type-3 optimality with the propagation premise supplied}
Fix any goal-preferring tie policy of \(A^\dagger\). On the domain
of solvable finite graphs with nonnegative costs and admissible
heuristics, no algorithm in \(\mathcal A_{\rm adm}\) strictly
uniformly dominates that policy. More precisely, a competitor that
skips one of its expanded states on some instance must, on another
admissible finite instance, expand a state it skips.

The comparison includes finite zero-cost edges. The proof's new terminal
edge can have zero cost; a comparison domain excluding all such edges
cannot simply reuse that construction. The competitor must satisfy the
correctness requirement, not merely be an arbitrary procedure.
\end{hbtheorem}

\begin{hbtheorem}{thm:p12-hsearch-no-type-two}{proved=annex:p12.hsearch-no-type-two}{No universally protected type-2 algorithm}
For every algorithm admissible on all solvable finite nonnegative-cost
instances, some other admissible algorithm has a proper-subset expansion
set on some such instance. Thus no fixed admissible algorithm satisfies
the pointwise type-2 test throughout this domain.
\end{hbtheorem}

\begin{hbtheorem}{thm:p12-hsearch-positive-type-three}{proved=annex:p12.hsearch-positive-type-three}{The strictly positive-cost comparison}
On the finite strictly positive-cost domain, use \(A^\dagger\) with
this additional tie convention: among minimum-barrier non-goals of
barrier \(m\), prefer a record with \(g<m\) to a record with
\(g=m\). Goals still have first priority; remaining ties are arbitrary.
Every fixed policy satisfying this convention is type-3 optimal against
algorithms admissible on all instances of this positive-cost domain.
No fixed admissible algorithm is type-2 optimal throughout that domain.

The same conclusions extend to the domain also allowing solvable locally
finite graphs with an existing positive uniform edge bound, when the
competitors are admissible there. No numerical lower bound common to all
instances is supplied as additional input: the finite comparison
instances may have arbitrarily small positive edges.
\end{hbtheorem}

\section{Generalized best-first barriers}
\label{sec:p12-hsearch-generalized}

\begin{hbdefinition}{def:p12-hsearch-general}{General scores and continuation-safe graph merging}
Let \(q(P)\) be any finite score of a finite path record. A general
best-first path search keeps separate records, selects a minimum-score
record, returns when its endpoint is a goal, and otherwise generates
all one-edge extensions. Ties prefer goals unless another convention is
explicitly stated. The initial non-goal is necessarily expanded first.

For a graph version with one record per state, require continuation-safe
order preservation of \(q\), as in
Definition~\ref{def:p12-hsearch-bestfirst-class}. Replace a record only
on a strict score improvement and reopen it if previously processed.
The cost \(C(P)\) assigned to a complete solution may be nonadditive;
stopping at a goal is part of the represented problem. No relationship
between \(q\) and \(C\) is assumed until stated.
In this generalized setting, \(C^*\) denotes the minimum of the
specified \(C(P)\) over allowed solution paths, whenever attained.
It agrees with the earlier additive \(h^*(s)\) when \(C(P)=g(P)\).
Assume each required score evaluation terminates. This section takes
\(s\) to be a non-goal; otherwise the empty solution is returned
immediately and no noninitial barrier is needed.

For a noninitial goal path put
\[
 B_q(P)=\max_{u\in P,\ u\text{ noninitial occurrence}}
                 q(P_{s,u}),\qquad
 \mu=\min_{P\text{ a goal path}}B_q(P).
\]
Assume the minimum is attained and finite. The forced initial selection
is excluded from this threshold. Including it gives the source's full
path maximum \(M=\max(q((s)),\mu)\).
\end{hbdefinition}

\begin{hbtheorem}{thm:p12-hsearch-general-barrier}{proved=annex:p12.hsearch-general-barrier}{Barrier, solution bound and strict expansion contour}
For the path-record or continuation-safe graph version:
\begin{enumerate}
\item After the initial expansion, until termination a current OPEN
 record has score at most \(\mu\); every subsequent selected record
 therefore has score at most \(\mu\).
\item If a solution record \(P_0\) is returned, then
 \(B_q(P_0)=\mu\). Some noninitial selection attains \(\mu\).
\item If goal scores are \(q(P)=\psi(C(P))\), with \(\psi\)
 strictly increasing, then \(\psi(C(P_0))\le\mu\). If an optimal
 solution path has its largest noninitial score at its goal, the returned
 solution is optimal.
\item Every non-goal state reachable by a path all of whose noninitial
 prefix scores are strictly below \(\mu\) is expanded before the
 first noninitial selection of score \(\mu\).
\end{enumerate}
The last two assertions concern terminating searches. The numerical
bound \(C(P_0)\le\psi^{-1}(\mu)\) requires that this inverse value
be defined. The weaker full-root bound with \(M\) also follows.
\end{hbtheorem}

\begin{hbtheorem}{thm:p12-hsearch-general-complete}{proved=annex:p12.hsearch-general-complete}{A general completeness condition}
Suppose the graph is locally finite, a solution exists, and along every
infinite rooted path-record chain the prefix scores are unbounded above.
Then the stated general best-first search terminates with a solution.
It suffices more generally that the prefix-closed tree of records
whose noninitial prefix scores are all at most \(\mu\) is finite.
\end{hbtheorem}

\begin{hblemma}{thm:p12-hsearch-general-closure}{proved=annex:p12.hsearch-general-closure}{What complete score propagation establishes}
In the continuation-safe graph version, suppose a finite revealed graph
has reached a score-relaxation fixed point: every revealed edge has been
relaxed after the last update of its tail's current record. Then the
current score at each reached state is the minimum score over all
start-rooted records in that revealed graph. Merely having generated
all edges of a record does not establish this fixed-point premise.
\end{hblemma}

\begin{hbexample}{def:p12-hsearch-saturated-cost}{A bounded score on an infinite positive-cost graph}
Let the initial state have a direct goal edge of cost 2 and an infinite
goal-free chain of unit-cost edges. Score a nonempty record by the
largest edge cost on it, and set the initial score to zero. Every chain
record has score 1, whereas the goal has score 2. General best-first
selection follows the chain forever. All edges cost at least 1, but
these scores are not unbounded above on the chain. Thus an edge-cost
bound that proves additive A* completeness cannot be transferred to
an arbitrary bounded evaluation score.
\end{hbexample}

\begin{hbtheorem}{thm:p12-hsearch-weighted}{proved=annex:p12.hsearch-weighted}{Weighted and nonlinear score certificates}
For additive costs, nonnegative \(h\), and a search preserving the
needed path frontier, the following assertions hold whenever it returns
a solution. Parts~1, 3 and~4 assume admissibility of \(h\); part~2 uses
its separately stated bound.
\begin{enumerate}
\item If \(q(P)=g(P)+h(v)(1+\varepsilon p(P))\), with
 \(\varepsilon\ge 0\) and \(p(P)\le 1\) along an optimal path,
 the returned cost is at most \((1+\varepsilon)C^*\).
\item If \(\beta\ge 0\) and \(h(v)\le\beta h^*(v)\) along an
 optimal path, the score
 \(q=(1-w)g+wh\), with \(0\le w<1\), is admissible when
 \(w\beta\le 1-w\). This includes
 \(w\le 1/(1+\beta)\), subject to the separate restriction \(w<1\).
\item A score \(q=H(g,h)\), nondecreasing in both arguments, is
 admissible if \(H(C,0)\) is strictly increasing in \(C\) and
 \(H(g,C-g)\le H(C,0)\) whenever \(0\le g\le C\).
 In particular \(H(g,h)=\sqrt{g^2+h^2}\) qualifies.
\item If \(\phi:[0,\infty)\to[0,\infty)\) is an increasing
 concave bijection with \(\phi(0)=0\), then
 \(H(g,h)=\phi(\phi^{-1}(g)+\phi^{-1}(h))\) qualifies.
\end{enumerate}
These are return-quality bounds, not independent termination theorems.
For a path-dependent score, invalid duplicate pruning cannot be excused
by one of these inequalities.
\end{hbtheorem}

\begin{hbexample}{def:p12-hsearch-nonlinear-failure}{A nonlinear combination that overestimates}
For \(H(g,h)=(\sqrt g+\sqrt h)^2\), use a branch
\(s\to r:1,r\to t:1\) and a direct goal edge \(s\to u:3\),
with \(h(s)=0,h(r)=1\) and zero goal heuristics. The score of
\(r\) is~4, while the direct goal's score is~3. Best-first selection
returns~3 although the optimum is~2. The input heuristic is admissible;
the nonlinear combination has destroyed its certificate.
\end{hbexample}

\begin{hbexample}{def:p12-hsearch-forced-root}{Why the initial score is not an expansion threshold}
In a finite tree give the initial node score~10, its non-goal child
\(a\) score~1, a goal sibling of \(a\) score~2, and a non-goal child
\(x\) of \(a\) score~9. The search selects the root, then \(a\),
then the goal, skipping \(x\). The full path minmax is~10, but the
effective threshold \(\mu\) is~2. A test using strictly less than~10
only after the root would wrongly force \(x\)'s expansion.
\end{hbexample}

\begin{hbtheorem}{thm:p12-hsearch-returned-path}{proved=annex:p12.hsearch-returned-path}{A branch test relative to the returned path}
Suppose the search returns the record
\(P=(s=v_0,v_1,\ldots,v_m=t)\), and put
\(L_i=\max_{i<j\le m}q(P_{s,v_j})\) for \(i<m\).
If a noninitial non-goal state \(x\) is expanded, there are a prefix
\(P_{s,v_i}\) and a continuation \(Q:v_i\leadsto x\) such that
every new prefix in the concatenated record has score at most \(L_i\).
Conversely, existence of such a continuation with every new prefix
score strictly less than \(L_i\) forces \(x\)'s expansion.
All scores use the complete start-rooted record, not the isolated branch
segment; the continuation contains no intervening goal.
In the necessary direction the prefix can be chosen as the common
initial segment of the returned record and a record at which \(x\)
was expanded.
\end{hbtheorem}

\section{An exact criterion on an ordered tree}
\label{sec:p12-hsearch-ordered-tree}

\begin{hbdefinition}{def:p12-hsearch-tree-convention}{Pure leftmost ties and side-subtree barriers}
Let \(T\) be a finite rooted tree with ordered children, and remove
all proper descendants of goals. In this section only, equal scores
are resolved purely from left to right, including ties involving goals.
This replaces, rather than supplements, the earlier goal-first tie rule.
Write \(q(u)\) for the unique path score at \(u\).
For a subtree rooted at \(r\), let \(m(r)\) be the minimum, over
its goal paths, of the largest score from \(r\) through the goal,
including both endpoints; set \(m(r)=+\infty\) if there is no goal.

Fix a target \(k\). For every noninitial \(n\) on the root-to-\(k\)
path, consider side subtrees branching at strict ancestors of \(n\).
Let \(\ell(n)\) be the minimum of their \(m(r)\)-values on the
left of the path, and \(r(n)\) the analogous minimum on the right.
The minimum of an empty collection is \(+\infty\). Side branches
at \(n\) itself are excluded: they are not generated before \(n\)
is selected.
\end{hbdefinition}

\begin{hbtheorem}{thm:p12-hsearch-tree-criterion}{proved=annex:p12.hsearch-tree-criterion}{Necessary and sufficient left/right inequalities}
The pure-leftmost search selects \(k\) if and only if, for every
noninitial \(n\) on its root path,
\[
                 q(n)<\ell(n),\qquad q(n)\le r(n).
\]
A selected non-goal is expanded; a selected goal terminates the search.
The initial node is selected without a side-subtree test. If the tree
has no goal, all its nodes are selected and both inequalities hold.
\end{hbtheorem}

\begin{hbexample}{def:p12-hsearch-tree-tie-boundary}{The goal-priority boundary}
Let the root have a left non-goal leaf and a right goal leaf, both of
score~1. The criterion selects the left leaf under pure-leftmost ties.
The earlier goal-first convention selects the right goal immediately
and skips the left leaf. Transferring the criterion without transferring
its tie convention changes its conclusion.
\end{hbexample}

\begin{hbtheorem}{thm:p12-hsearch-no-type-zero}{proved=annex:p12.hsearch-no-type-zero}{Why universal tie-independent containment can fail}
On the finite nonnegative-cost consistent domain, no single successful
search algorithm has an expansion set contained in every admissible
competitor's on every instance. The obstruction already uses
uniform-cost competitors, so it also applies when competitors are
restricted to the normalized path-score class or to the costwise-compatible
class. Nonpathological-domain containment is not contradicted.
\end{hbtheorem}

An expansion theorem remains relative to its algorithm class and cost
measure. Re-expansion counts, heuristic-evaluation costs, priority-queue
work, memory and elapsed time are not determined by containment of
distinct expanded states. A heuristic obtained by a proved relaxation
offers a useful certificate; an empirical claim that most heuristics
people invent satisfy it would require separate evidence.

\begin{hbcorollary}{thm:p12-hsearch-positive-domain-extension}{proved=annex:p12.hsearch-positive-domain-extension}{The source's locally finite positive-cost domain}
The strict-contour, consistent-containment, nonpathological-trade,
matched-tie A** containment and normalized path-score conclusions extend
to solvable locally finite graphs whose edge costs have an existing
positive uniform lower bound. For each conclusion, use comparison
algorithms satisfying its stated correctness requirement on that entire
domain. The weak-compatibility counterexample, the pending-A** domination
counterexample and the separate costwise theorem extend in the same sense;
the latter two compatibility requirements retain their stated quantifiers
when advice overestimates.

For a normalized two-input path score in this positive-cost domain,
\(H(x,y)\le x+y\) is forced for \(x>0,y\ge 0\). No restriction
on the compulsory root's score is inferred by treating it as a noninitial
state of zero prefix cost. Failure of universal containment against
admissible and normalized path-score competitors already occurs on
finite strictly positive-cost consistent graphs.
\end{hbcorollary}
